\documentclass[10pt,a4paper]{amsart}
\usepackage[margin=19mm]{geometry}
\usepackage{amsmath,amssymb,mathtools}
\usepackage[scr=rsfs,cal=euler]{mathalfa}
\usepackage{xfrac}
\usepackage[unicode,hidelinks,bookmarks=false]{hyperref}
\newcommand{\ie}{i.e.\ }
\newcommand{\cf}{cf.\ }
\newcommand{\resp}{respectively\ }
\newcommand{\Subsection}{\subsection}
\providecommand{\nobreakdash}{\nobreak\hbox{-}}
\setlength{\emergencystretch}{2em}
\multlinegap0pt
% ---------------------------------------------------------------------------
% Editorial AMS wrapper prelude. The theorem environments and the mathematical macros below are
% transcribed from the paper's own preamble (cached-originals/source0.tex lines 29--232); the AMS amsart
% class, the named format packages of the pinned TeX Live runtime and this editorial formatting are not
% part of the author's text. No mathematics is changed.
% ---------------------------------------------------------------------------
\newtheorem{Th}{Theorem}[section]
\newtheorem{Prop}[Th]{Proposition}
\newtheorem{Lem}[Th]{Lemma}
\newtheorem{Cor}[Th]{Corollary}
\newtheorem{Def}[Th]{Definition}
\newtheorem{Rem}[Th]{Remark}       
\newtheorem{Ex}[Th]{Example}
\newtheorem{Ass}[Th]{Assumption}
\newtheorem{Que}[Th]{Question}

\newenvironment{altproof}[1]
{\noindent%\addvspace{0.3cm}
{\em Proof of {#1}}.}
{\nopagebreak\mbox{}\hfill $\Box$\par\addvspace{0.5cm}}

\newcommand{\wt}{\widetilde}
\newcommand{\mf}{\mathfrak}
   \newcommand{\vp}{\varphi}
   \newcommand{\Vp}{\Phi}
   \newcommand{\eps}{\varepsilon}
\newcommand{\ov}{\overline}
   \def\li{\mathop{\mathrm{liminf}\,}}
   \def\div{\mathop{\mathrm{div}}}
   \def\ls{\mathop{\mathrm{limsup}\,}}
   \def\Li{\mathop{\mathrm{Liminf}\,}}
   \def\Ls{\mathop{\mathrm{Limsup}\,}}
   \def\deg{\mathrm{deg}}
     \def\span{\mathrm{span}}
   \def\ind{\mathrm{ind}}
   \def\supp{\mathrm{supp}}
   \def\lan{\langle}
   \def\ran{\rangle}
   \def\conv{\mathrm{conv}\,}
   \def\id{\mathrm{id}}
    \def\O{\mathrm{O}}
   \def\Q{\mathbb{Q}}
   \def\Z{\mathbb{Z}}
   \def\d{\diamond}
   \def\bd{\mathrm{bd}\,}
   \def\multi{\multimap}
   \def\multito{\mapstochar\multimap}
   \def\N{\mathbb{N}}
   \def\R{\mathbb{R}}
   \def\Gr{\mathrm{Gr}\,}
   \def\Ind{\mathrm{Ind}}
   \def\ind{\mathrm{ind}}
   \def\inr{\mathrm{int}}
   \def\curl{\mathrm{curl}}
  \def\dim{\mathrm{dim}}
   \def\defdim{\mathrm{defdim}}
   \def\diam{\mathrm{diam}}
   \def\Gr{\mathrm{Gr}}
   \def\st{\mathrm{st}}
   \def\cl{\mathrm{cl\,}}
   \def\Fix{\mathrm{Fix}}
   \def\F{{\cal F}}

   \def\P{\mathcal P} %P family
   \def\Uc{\mathscr{U}} %U cover
   \def\U{\mathcal{U}} %U neighborhood of graph
   \def\Vc{\mathscr{V}}
   \def\V{\mathcal{V}}
   \def\E{\mathcal{E}}
   \def\J{\mathcal{J}}
   \def\G{\mathcal{G}}
   \def\Wc{\mathscr{W}}
   \def\W{\mathcal{W}}
   \def\D{\mathcal{D}}
   \def\C{\mathbb{C}}
   \def\T{\mathcal{T}}
   \def\Dc{\mathscr{Dc}}
   \def\A{\mathcal{A}}
     \def\M{\mathcal{M}}
     \def\tJ{\wt{\cJ}}
    
    
    \newcommand{\DF}{\cD_{\cF}}
    \newcommand{\UU}{u}
    \newcommand{\VV}{\mathbf{V}}
    \newcommand{\BB}{\mathbf{B}}
    \newcommand{\AAa}{\mathbf{A}}
    \newcommand{\mbS}{\mathbb{S}}
    \newcommand{\SO}{\cS\cO}
    

\def\rh{\rightharpoonup}
\def\irn{\int_{\rn}}
\def\rn{\mathbb{R}^N}
\newcommand{\cA}{{\mathcal A}}
\newcommand{\cB}{{\mathcal B}}
\newcommand{\cC}{{\mathcal C}}
\newcommand{\cD}{{\mathcal D}}
\newcommand{\cDp}{\cD^{1,p}(\R^3,\R^3)}
\newcommand{\cE}{{\mathcal E}}
\newcommand{\cF}{{\mathcal F}}
\newcommand{\cG}{{\mathcal G}}
\newcommand{\cH}{{\mathcal H}}
\newcommand{\cI}{{\mathcal I}}
\newcommand{\cJ}{{\mathcal J}}
\newcommand{\cK}{{\mathcal K}}
\newcommand{\cL}{{\mathcal L}}
\newcommand{\cM}{{\mathcal M}}
\newcommand{\cN}{{\mathcal N}}
\newcommand{\cO}{{\mathcal O}}
\newcommand{\cP}{{\mathcal P}}
\newcommand{\cQ}{{\mathcal Q}}
\newcommand{\cR}{{\mathcal R}}
\newcommand{\cS}{{\mathcal S}}
\newcommand{\cT}{{\mathcal T}}
\newcommand{\cU}{{\mathcal U}}
\newcommand{\cV}{{\mathcal V}}
\newcommand{\cW}{{\mathcal W}}
\newcommand{\cX}{{\mathcal X}}
\newcommand{\cY}{{\mathcal Y}}
\newcommand{\cZ}{{\mathcal Z}}

\newcommand{\fA}{{\mathfrak A}}
\newcommand{\fB}{{\mathfrak B}}
\newcommand{\fC}{{\mathfrak C}}
\newcommand{\fD}{{\mathfrak D}}
\newcommand{\fE}{{\mathfrak E}}
\newcommand{\fF}{{\mathfrak F}}
\newcommand{\fG}{{\mathfrak G}}
\newcommand{\fH}{{\mathfrak H}}
\newcommand{\fI}{{\mathfrak I}}
\newcommand{\fJ}{{\mathfrak J}}
\newcommand{\fK}{{\mathfrak K}}
\newcommand{\fL}{{\mathfrak L}}
\newcommand{\fM}{{\mathfrak M}}
\newcommand{\fN}{{\mathfrak N}}
\newcommand{\fO}{{\mathfrak O}}
\newcommand{\fP}{{\mathfrak P}}
\newcommand{\fQ}{{\mathfrak Q}}
\newcommand{\fR}{{\mathfrak R}}
\newcommand{\fS}{{\mathfrak S}}
\newcommand{\fT}{{\mathfrak T}}
\newcommand{\fU}{{\mathfrak U}}
\newcommand{\fV}{{\mathfrak V}}
\newcommand{\fW}{{\mathfrak W}}
\newcommand{\fX}{{\mathfrak X}}
\newcommand{\fY}{{\mathfrak Y}}
\newcommand{\fZ}{{\mathfrak Z}}
\renewcommand{\dim}{{\rm dim}\,}
\newcommand{\vphi}{\varphi}
\newcommand{\al}{\alpha}
\newcommand{\be}{\beta}
\newcommand{\ga}{\gamma}
\newcommand{\de}{\delta}
\newcommand{\la}{\lambda}
\newcommand{\De}{\Delta}
\newcommand{\Ga}{\Gamma}
\newcommand{\Om}{\Omega}

\def\r{\mathbb{R}}
\def\r3{\mathbb{R}^3}
\def\Wc{\cD^{1,p}(\curl;\R^3)}
\def\Wco{\cD^{1,p}_{\cO}(\curl;\R^3)}
\def\rn{\mathbb{R}^N}
\def\rm{\mathbb{R}^{M}}
\def\z{\mathbb{Z}}
\def\zn{\mathbb{Z}^N}
\def\n{\mathbb{N}}
\def\ch{\mathcal{H}}
\def\eps{\varepsilon}
\def\rh{\rightharpoonup}
\def\irn{\int_{\rn}}
\def\irm{\int_{\rm}}
\def\io{\int_{\Omega}}
\def\iom{\int_{\omega}}
\def\ir3{\int_{\r3}}
\def\vp{\varphi}
\def\wt{\widetilde}
\def\wh{\widehat}
\def\ol{\overline}
\def\la{\langle}
\def\ra{\rangle}
\def\lam{\lambda}
\def\lam1{\lambda_1}
\def\donetwo{\mathcal{D}^{1,2}}
\def\d1p{\mathcal{D}^{1,p}}
\def\cm{\mathcal{M}}
\def\meas{\text{meas}}

\def\curlop{\nabla\times}
\newcommand{\weakto}{\rightharpoonup}
\newcommand{\pa}{\partial}
\def\id{\mathrm{id}}
\def\ran{\mathrm{range}}
\newcommand{\tX}{\widetilde{X}}
\newcommand{\tu}{\widetilde{u}}
\newcommand{\tv}{\widetilde{v}}
\newcommand{\tcV}{\widetilde{\cV}}


\newcommand{\wti}{\widetilde}
\newcommand{\cTto}{\stackrel{\cT}{\longrightarrow}}

\numberwithin{equation}{section}
\newcommand{\open}{\quad{\bf To be checked/done!}\quad}


\begin{document}
\title{Optimizers in Sobolev-curl inequalities}
\author{Jaros\l aw Mederski$^{*}$ and Andrzej Szulkin}
\date{Source: arXiv:2511.01432v2; distributed under CC BY 4.0}
\maketitle
\noindent\emph{Source, version and licence.} \emph{Optimizers in Sobolev-curl inequalities}, by Jaros\l aw Mederski$^{*}$ and Andrzej Szulkin. The source of this editable unit is the e-print of \textbf{version v2} of arXiv:2511.01432, https://arxiv.org/abs/2511.01432v2, whose material is distributed under the Creative Commons Attribution 4.0 International licence, http://creativecommons.org/licenses/by/4.0/, as recorded in the arXiv licence metadata for this paper and shown on that abstract page. The whole of the body below is version v2, the newest version of the paper. Full rights notice: licenses/arxiv-2511.01432-CC-BY-4.0.txt.\par\smallskip
\noindent\textit{Editorial scope.} Counted unit: original Lemma 2.1 of the article, the statement carried by the source environment \texttt{Lem} with source label \texttt{defof}, cached-originals/source0.tex lines 561--573, printed as \textbf{Lemma 2.1} exactly as in the article. It is delivered together with the author's own complete proof as the single primary proof body, source0.tex lines 577--626 (50 lines): the \texttt{proof} environment that immediately follows the statement and establishes the closedness of both spaces and the norm comparison. No proof marker is added and no mathematics is written, repaired or re-derived; the author's notation and the printed slips of the source are kept literal. Necessary complete context is retained uncounted, in source order: source0.tex lines 337--339 and 365--368, the definitions and identities of the spaces $\cD^{1,p}$ and of the normals used in the statement, which the closure of the counted proof requires; and lines 430--438, Theorem 1.1, the first main theorem of the article whose proof uses the counted lemma. The article's own numbering is reproduced by explicit counter settings: the counter \texttt{Th}, declared by \texttt{\textbackslash newtheorem\{Th\}\{Theorem\}[section]} and shared by Prop, Lem, Cor, Def, Rem, Ex, Ass and Que, is set before each retained passage, so the counted statement prints as Lemma 2.1 and the main theorem prints as Theorem 1.1, exactly as in the article. Every \texttt{\textbackslash ref} and \texttt{\textbackslash eqref} inside every retained passage resolves inside this delivered file, and the three citations that occur in the retained passages, \texttt{Iwaniec}, \texttt{Le} and \texttt{MSchSz}, are delivered as the authors' own bibliography entries transcribed verbatim from the article's own \texttt{thebibliography}, keeping the authors' own printed numbers. The retained passages contain no figure environment and no graphics-inclusion command, so no artwork is redistributed. The source's own class is AMS \texttt{amsart} (\texttt{\textbackslash documentclass[12pt,leqno,twoside]\{amsart\}}); the e-print ships no class or style file, so no publisher class is shipped, used or redistributed, and the wrapper is the same AMS \texttt{amsart} class with the paper's own macros and theorem declarations transcribed from its preamble. Every declared body of this unit is an exact, unmodified span of source0.tex: no body is a bounded passage comparison, \texttt{omit} was not used anywhere, and consequently no fidelity receipt is asserted in delivery-evidence.json. The complete concatenated original source is bundled as sources/188318/source0.tex. Omitted from this editable unit and therefore absent from it are source0.tex lines 1--336; 340--364; 369--429; 439--560; 574--576; 627--1453; 1462--1473; 1476--1513, that is the article's own preamble, title block and abstract, the introduction with the remaining main theorems, the sections on the concentration-compactness machinery and on the symmetry of optimizers with their lemmas and proofs, and the remaining bibliography entries, all of which remain present in the bundled source but are not delivered here. The AMS \texttt{amsart} wrapper, the named format packages of the pinned TeX Live runtime and this editorial source, licence and scope paragraph are editorial formatting only.\par\smallskip
\setcounter{section}{1}
\setcounter{equation}{3}
\setcounter{Th}{0}
\begin{equation}\label{eq:neq}
\int_{\r3}|\curlop u|^p\, dx\geq S_{p,\curl}\inf_{w\in \cW}\Big(\int_{\r3}|u+w|^{p^*}\,dx\Big)^{\frac{3-p}{3}}
\end{equation}

\setcounter{section}{1}
\setcounter{equation}{5}
\setcounter{Th}{0}
\begin{equation}
\label{eq}
\curlop (|\curlop u|^{p-2}\curlop u)= |u|^{p^*-2}u\quad \hbox{in }\R^3.
\end{equation}

\setcounter{section}{1}
\setcounter{equation}{11}
\setcounter{Th}{0}
\begin{Th}\label{Th:main1}
(a) $S_{p,\curl} > S_p\cdot\cH_p$. \\
(b)  If $(u_n)\subset \cN$ is a minimizing sequence for $J$, then there are $(s_n)\subset (0,\infty)$ and $(y_n)\subset \R^3$ such that, passing to a subsequence, 
$$
s_n^{3/p^*}u_n(s_n\cdot+y_n)\to u
$$ 
where $u$ is a minimizer for $J$ on $\cN$.\\
(c) $\inf_{\cN}J=\frac13S_{p,\curl}^{3/p}$ is attained,  $u$ is a ground state solution to \eqref{eq} and equality holds in \eqref{eq:neq} for this $u$. If $u\in\Wc\setminus\cW$ satisfies equality in \eqref{eq:neq}, then there are unique $t>0$ and $w\in\cW$ such that $t(u+w)\in\cN$ and $J(t(u+w))=\inf_{\cN} J$.
\end{Th}

\setcounter{section}{2}
\setcounter{equation}{0}
\setcounter{Th}{0}
\begin{Lem}\label{defof} $\V$ and $\cW$ are closed subspaces of $\Wc$  and
	\begin{equation}\label{HelmholzDec}
	\Wc=\V\oplus \W
	\end{equation}
where $$\cV:=\Big\{v\in \Wc: \int_{\R^3}\langle v,\nabla \varphi\rangle\,dx=0
\text{ for every $\varphi\in \cC^\infty_0(\R^3)$} \Big\}$$
and
$$
\cW:=\Big\{w\in \Wc: \int_{\R^3}\langle w,\curlop \varphi\rangle\,dx=0
\text{ for every $\varphi\in \cC^\infty_0(\R^3,\R^3)$} \Big\}.
$$
Moreover, $\cV\subset\cD^{1,p}(\r3,\r3)$ and the norms $|\nabla \cdot|_p$ and $|\curlop \cdot|_p$ are equivalent in $\V$. 
\end{Lem}

\setcounter{section}{2}
\setcounter{equation}{1}
\setcounter{Th}{1}
\begin{proof} 
We follow similar arguments as in \cite[Lemma 2.4]{MSchSz} for $p=2$; however, since in general for divergence-free $u$, $|\curlop u|_p\neq |\nabla u|_p$ if $p\neq 2$, we have to modify the proof.
 
 Firstly, we easily check that $\V$ and $\cW$ are closed subspaces of $\Wc$.
	 
Now, take any $u\in \Wc$ and $u_n\in\cC_0^{\infty}(\R^3,\R^3)$ such that $u_n\to u$ in $\Wc$. Let $\vp_n\in\cC^{\infty}(\R^3)$ be the Newtonian potential of $\div(u_n)$, i.e. $\vp_n$ solves $\Delta \vp_n = \div(u_n)$.  Since $u_n\in\cC_0^{\infty}(\R^3,\R^3)$, then by \cite[Proposition 1]{Iwaniec}, $\nabla \vp_n\in L^r(\R^3,\R^3)$ for every $r\in(1,\infty)$. In particular, $\nabla\vp_n\in L^{p^*}(\R^3,\R^3)$. 
Hence
$$	v_n:=u_n-\nabla \vp_n\in  L^{p^*}(\R^3,\R^3).
	$$ 
	Note also that  $\curlop v_n = \curlop u_n$  and $\div(v_n)=0$ pointwise. Let $v_n=(v_n^1,v_n^2,v_n^3)$ and let $e_i$ denote the $i$-th element of the standard basis in $\r3$ (so $v_n^i = \langle v_n,e_i\rangle$). Since
$$-\Delta  v_n=\curlop \curlop v_n,$$
employing the vector calculus identity
\[
\div(A\times B) = \langle \curlop A, B\rangle - \langle A,\curlop B\rangle
\] 
with $A = \curlop v_n$ and $B=e_i$, we obtain
\[
-\Delta v_n^i = \langle \curlop(\curlop v_n), e_i\rangle = \div((\curlop v_n)\times e_i) =: \div(f_n^i), \quad i=1,2,3. 
\]
Note that $\curlop v_n = \curlop u_n\in L^2(\R^3,\R^3)\cap L^p(\R^3,\R^3)$ and
$f_n^i\in L^2(\R^3,\R^3)\cap L^p(\R^3,\R^3)$. Again  by \cite[Proposition 1]{Iwaniec}, 
	$$|\nabla v_n^i|_p\leq C_p|f_n^i|_p \quad\hbox{for }i=1,2,3$$
where $C_p>0$ is a constant.
Hence 
\begin{equation}\label{eqes3C}
|\nabla v_n|_p\leq D_p |\curlop v_n|_p = D_p |\curlop u_n|_p,
\end{equation}
for a constant $D_p>0$ depending on $C_p$.
Similarly we obtain that for $m,n\ge 1$,
	$$|\nabla(v_n-v_m)|_{p}\leq D_p|\curlop(v_n-v_m)|_{p}=D_p|\curlop(u_n-u_m)|_{p}\leq D_p \|u_n-u_m\|.$$
	Thus $(v_n)$ is a Cauchy sequence in $\cDp$. Let $v:=\lim_{n\to\infty}v_n$ in $\cDp$. Then
	$$\int_{\R^3}\langle v,\nabla \vp\rangle\,dx=\lim_{n\to\infty}\int_{\R^3}\langle v_n,\nabla \vp\rangle\,dx=0$$
	for any $\vp\in\cC_0^{\infty}(\R^3)$, hence $\div(v) =0$ and $v\in\cV$.
	Moreover, as $|\curlop v|^2 \le 2|\nabla v|^2$, we have 
	\begin{equation} \label{ineq5}
	|\curlop(v_n-v)|_{p}\leq 2^{1/2}|\nabla(v_n-v)|_{p}\to 0,
	\end{equation}
	so $v_n\to v$ in  $\Wc$ and
	$\nabla \vp_n=u_n-v_n\to u-v$ in $\Wc$. Since $\nabla\vp_n\in\cW$ and $\cW$ is closed,
	$u-v\in \cW$ and we get the decomposition
	$$u=v+(u-v)\in \cV+\cW.$$
	\indent Now take $v\in \cV\cap\W$ and since $\curlop v =0$,
	 by \cite[Lemma 1.1(i)]{Le}, $v=\nabla\xi$ for some $\xi\in W^{1,{p^*}}_{loc}(\R^3)$. Since $\div(v) = 0$, $\xi$ is harmonic and so is $v$. Since $v\in L^{p^*}(\R^3,\R^3)$, by the mean-value formula we infer that $v=0$, so \eqref{HelmholzDec} holds.

From \eqref{eqes3C} with $v$ replacing $v_n$ and \eqref{ineq5} with $v$ replacing $v_n-v$ we see that
\begin{equation} \label{doubleineq}
2^{-1/2}|\curlop v|_p \le  |\nabla v|_p\leq D_p |\curlop v|_p,
\end{equation}
i.e. the norms $|\nabla\cdot|_p$ and $|\curlop\cdot|_p$ are equivalent in $\cV$.
\end{proof}

\begin{thebibliography}{99}
\bibitem[18]{Iwaniec} T. Iwaniec: {\em Projections onto gradient fields and $L^p$-estimates for degenerated elliptic operators}, Studia Math. {\bf 75} (1983), no. 3, 293--312. 
\bibitem[19]{Le} H. Leinfelder: \emph{Gauge invariance of Schr\"odinger operators and related spectral properties}, J. Operator Theory 9 (1983), 163--179. 
\bibitem[25]{MSchSz} J. Mederski, J. Schino, A. Szulkin: {\em Multiple solutions to a nonlinear curl-curl problem in $\R^3$}, Arch. Rational Mech. Anal. 236 (1) (2020), 253--288. 
\end{thebibliography}
\end{document}
